\infographicchapter{Dual-Energy X-Ray Absorptiometry}{dexa.png}
\label{ch:infographic-dexa}

\section{Overview}

Dual-energy X-ray absorptiometry (often called a DEXA or DXA scan) is a quick, gentle, ultra-low-dose X-ray test and the gold standard for measuring bone mineral density---how strong and dense your bones are. Unlike a regular X-ray that uses one energy level, a DXA scanner sends two very gentle X-ray beams of different strengths through your body. Because soft tissue (like muscle and fat) and hard bone absorb those two beams differently, a computer cleanly subtracts the soft tissue away and calculates the exact amount of mineral packed into your bones (measured in grams per square centimeter, or g/cm\textsuperscript{2}).

A standard DXA exam checks the bones most likely to break if bone strength drops: your lower back (lumbar spine) and hip (and sometimes your forearm if the spine or hip cannot be scanned). Whole-body DXA scans can also measure body fat and lean muscle mass. Your results come with two easy-to-track scores. A T-score compares your bone density to that of a healthy young adult: for postmenopausal women and men aged 50 or older, a T-score of -1.0 or higher is normal, between -1.0 and -2.5 means low bone mass (osteopenia), and -2.5 or lower means osteoporosis. A Z-score compares your bones to other people of your own age and sex, which is used for younger adults and children (a Z-score of -2.0 or lower means bone density is lower than expected for age). Doctors use DXA to screen for osteoporosis, estimate your 10-year fracture risk (often using a calculator called FRAX), and check whether bone-strengthening medicines are working.

DXA is fast, completely painless, open-air (no tunnel), and uses an ultra-low radiation dose---similar to just 1--2 days of natural outdoor background exposure and far less than a single chest X-ray---making repeat checkups very safe. Its main limits are that it measures bone in two flat dimensions, so arthritis bone spurs, old spine compression fractures, or calcium in nearby blood vessels can make spine numbers look falsely higher than they really are, and standard DXA does not directly show fine internal bone structure.

\section{Underlying Physical Principles}
\begin{itemize}
\item \textbf{Two Gentle X-Ray Energy Levels:} A DXA machine produces two distinct low-dose X-ray beams---one lower energy and one higher energy---either by switching the X-ray tube voltage rapidly or by using a special metal filter.
\item \textbf{How Bone and Soft Tissue Behave Differently:} Hard, calcium-rich bone absorbs low-energy X-rays much more strongly than high-energy X-rays, whereas overlying soft tissue (fat and muscle) absorbs both energies much more evenly.
\item \textbf{Cleanly Subtracting Soft Tissue:} By comparing how much of each X-ray beam passes through your body at every point, the computer mathematically cancels out overlying fat and muscle, leaving an accurate measurement of bone mineral content alone.
\item \textbf{Calculating Bone Density and T-Scores:} A fan-shaped X-ray beam and detector glide over your lower spine, hip, or forearm, dividing total bone mineral content by the scanned bone area to find areal bone mineral density (g/cm\textsuperscript{2}). The software then compares that number to healthy reference charts to generate your T-score and Z-score.
\end{itemize}

\section{Imaging Mechanism}
\begin{enumerate}
\item \textbf{Daily Accuracy Check and Positioning:} Each morning, staff test the scanner on a standardized bone-like block (a phantom) to ensure pinpoint accuracy. You then lie flat on your back on an open padded table, using a foam block under your lower legs for the spine scan or a gentle foot brace that turns your leg slightly inward for the hip scan.
\item \textbf{Smooth Overhead Scanning Pass:} While you lie still, an X-ray source beneath the table and a detector arm overhead glide smoothly together over your lower back, hip, or forearm.
\item \textbf{Counting Dual-Energy Signals:} Solid-state detectors in the overhead arm count the low-energy and high-energy X-rays passing through each millimeter of the scanned area.
\item \textbf{Subtracting Soft Tissue and Outlining Bone:} Densitometry software cancels out soft tissue pixel by pixel, outlines the lower spine bones (L1--L4) or hip zones automatically, and calculates bone mineral density.
\item \textbf{Generating Your Score Report:} The computer compares your numbers to national reference data (such as NHANES III) and displays verification images alongside your T-scores, Z-scores, trend graphs from past scans, and optional spine-fracture or bone-texture scores.
\end{enumerate}

\section{Key Technical Parameters}
\noindent
\begin{tabularx}{\textwidth}{@{}l X@{}}
\toprule
\textbf{Parameter} & \textbf{Typical Value / Description} \\
\midrule
Physical Principle & Dual-energy X-ray scanning that subtracts soft tissue to measure bone mineral density (g/cm\textsuperscript{2}) \\
Energy / Signal Source & Dual-energy X-ray tube paired with a flat solid-state linear detector arm \\
Typical Frequency / Energy & Two gentle X-ray energy peaks: about 40--70 keV (low) and 70--140 keV (high) \\
Spatial Resolution & About 0.5--1.0 mm pixel spacing (built for precise bone-density numbers rather than fine fracture X-rays) \\
Typical Scan Duration & 10--20 minutes total visit; only 1--3 minutes per scanned area (spine, hip, or forearm) \\
Contrast Agent & None (recent barium, CT contrast dye, or nuclear tracers interfere with the scan) \\
Ionizing Radiation & Yes; ultra-low effective dose of about 1--10 microsieverts (similar to 1--2 days of natural background radiation) \\
\bottomrule
\end{tabularx}

\section{Common Symptoms That Prompt This Scan}
\begin{itemize}
\item \textbf{Often Silent (No Symptoms at All):} Bone thinning usually causes no pain or early warning signs, so most DEXA scans are prompted by age or risk-factor screening rather than symptoms.
\item \textbf{Broken Bone From a Minor Fall:} A wrist, hip, spine, or shoulder fracture after a simple slip from standing height or an everyday bump.
\item \textbf{Noticeable Height Loss or Stooped Posture:} Losing half an inch (about 1.3~cm) or more within a single year, losing more than 1.5~inches (about 4~cm) of total adult height, or developing a gradually curving upper back.
\item \textbf{New Mid- or Lower-Back Pain in Older Adults:} Sudden or persistent localized back pain that may signal a hidden spinal compression fracture.
\end{itemize}

\section{Clinical Usages}

\subsection{Primary Usages}
\begin{itemize}
\item \textbf{Osteoporosis Screening:} The gold-standard test for diagnosing osteoporosis and osteopenia in women aged 65 or older, men aged 70 or older, and younger postmenopausal women or men aged 50--69 who have risk factors such as prior fractures, family history, very low body weight, rheumatoid arthritis, diabetes, long-term tobacco or alcohol use, or low calcium and vitamin D intake.
\item \textbf{Estimating Future Fracture Risk:} Supplies hip bone density measurements for risk calculators like FRAX, which estimate your 10-year chance of a hip or major bone fracture to help decide whether medication is needed.
\item \textbf{Checking for Medication- or Illness-Related Bone Loss:} Monitors bone strength in people taking long-term steroid medicines (glucocorticoids) or certain hormone-blocking cancer treatments, as well as people with thyroid/parathyroid conditions, malabsorption, or chronic kidney disease.
\item \textbf{Tracking Treatment Results:} Measures changes in spine and hip bone density every 1--2 years to confirm that osteoporosis medicines are stabilizing or rebuilding bone.
\end{itemize}

\subsection{Secondary Usages}
\begin{itemize}
\item \textbf{Spotting Silent Spine Fractures (VFA):} A quick side-view spine scan (Vertebral Fracture Assessment) detects painless, hidden spine compression fractures that strongly raise future fracture risk.
\item \textbf{Checking Internal Bone Texture (TBS):} Special software called Trabecular Bone Score evaluates the fine texture of your lower spine images to add extra insight into bone strength beyond standard density numbers.
\item \textbf{Measuring Body Composition:} Whole-body DXA accurately measures total and regional body fat, deep abdominal fat, and lean muscle mass for evaluating muscle loss (sarcopenia), weight management, or athletic training.
\item \textbf{Checking Bone Health in Children and Young Adults:} Uses age- and sex-matched Z-scores (adjusted for height when needed) to monitor bone fragility in children with brittle-bone conditions, neuromuscular disorders, or chronic inflammatory illness.
\end{itemize}

\section{Advantages}
\begin{itemize}
\item \textbf{Trusted Gold Standard:} Provides proven T-scores and Z-scores used worldwide by the World Health Organization and FRAX fracture calculators.
\item \textbf{Ultra-Low Radiation Dose:} Uses only a tiny fraction of the radiation of a standard chest X-ray---about equal to 1--2 days of everyday background exposure---making repeat monitoring very safe.
\item \textbf{Highly Accurate and Repeatable:} Very consistent measurements allow doctors to track subtle improvements or losses in bone density reliably over 1--2 year intervals.
\item \textbf{Fast, Painless, and Open-Air:} Takes just minutes on an open padded table with no enclosed tunnel, no injections, and no discomfort.
\item \textbf{Helpful Add-On Information:} Can be paired in the same visit with hidden spine-fracture screening (VFA), bone-texture scoring (TBS), or full-body fat and muscle analysis.
\end{itemize}

\section{Disadvantages}
\begin{itemize}
\item \textbf{Sensitive to Bone Size:} Because DXA measures bone density over a flat two-dimensional area rather than a full 3D volume, people with very small frames can sometimes appear to have slightly lower bone density than they truly do.
\item \textbf{Arthritis and Calcium Can Skew Spine Numbers:} Spinal arthritis bone spurs, old compressed vertebrae, calcium deposits in the main belly artery (aorta), or metal surgical hardware can make lower-back bone density look falsely high.
\item \textbf{Does Not Directly Image Fine 3D Structure:} Standard DXA cannot separate the hard outer shell of bone from the spongy inner core on its own without add-on tools like TBS or specialized CT scans.
\item \textbf{Same-Machine Consistency Matters:} Different DXA scanner brands calibrate slightly differently, so accurate tracking over time works best when you return to the same facility and machine with consistent positioning.
\item \textbf{Pregnancy and Interference Limits:} Routine DXA is avoided during pregnancy due to X-rays; recent barium, CT contrast, or nuclear scans can throw off the numbers; and very high body weight may exceed some scanner table limits.
\end{itemize}

\section{Contraindications}

\subsection{Absolute}
\begin{itemize}
\item \textbf{Pregnancy:} Routine DXA is not performed during pregnancy because it uses X-rays and bone density testing can safely wait until after delivery.
\end{itemize}

\subsection{Relative}
\begin{itemize}
\item \textbf{Recent Contrast Dye or Nuclear Scans:} Having a barium X-ray, a CT scan with oral or IV contrast dye, or a nuclear medicine scan within the past 7--14 days can interfere with DXA measurements, so the bone scan should be rescheduled once those clear.
\item \textbf{Table Weight or Lying-Flat Limits:} Body weight above the scanner table's limit (typically 350--500 pounds depending on the model) or trouble lying flat for a few minutes may prevent spine and hip scanning, though a seated forearm DXA can often be done instead.
\item \textbf{Metal Hardware in Both Hips or Heavy Spine Surgery:} Metal replacements in both hips or extensive spinal rods and screws block accurate central measurements, so the technologist scans the non-dominant forearm (33\% radius) instead.
\end{itemize}

\section{Patient Safety Guidelines}
\begin{itemize}
\item \textbf{Pregnancy Check and Minimal Radiation:} Even though the X-ray dose is ultra-low, staff ask all women who could be pregnant about pregnancy status and postpone routine DXA scans during pregnancy.
\item \textbf{Checking for Interfering Recent Scans:} Staff confirm you have not had a barium study, contrast CT, or nuclear scan in the past 7--14 days so you do not get inaccurate numbers or need a repeat scan.
\item \textbf{Safe Table Support and Fall Prevention:} Staff verify the table weight limit, help frail or elderly patients safely onto and off the table, and position tender hips or backs gently.
\item \textbf{Keeping Measurements Clean and Consistent:} Staff ask you to remove metal fasteners and jewelry, pause calcium pills on exam day, and ideally scan you on the same calibrated machine for follow-up visits.
\item \textbf{Daily Calibration Checks:} Clinics run daily phantom tests and keep detailed quality logs so your bone-density numbers remain completely trustworthy for treatment decisions.
\end{itemize}

\section{Prerequisite Tests}
\begin{itemize}
\item \textbf{No Blood Tests Needed to Have the Scan:} You do not need fasting or blood tests just to get a DXA scan. However, if your scan reveals osteopenia or osteoporosis, your doctor will usually order simple blood tests (calcium, vitamin D, kidney function, blood count, and thyroid hormone) to check for hidden causes of bone loss.
\item \textbf{Height, Weight, and Health History:} Staff measure your exact standing height and weight and note your menopausal status, past fractures, steroid use, and risk factors to calculate accurate T-scores, Z-scores, and FRAX risk.
\item \textbf{Pregnancy Check:} Women of childbearing age are asked about possible pregnancy (with a quick pregnancy test if uncertain) so routine scans can be postponed during pregnancy.
\item \textbf{Reviewing Recent Scans and Metal Hardware:} Staff check for any recent contrast or nuclear scans within the last 7--14 days and note any hip replacements, spinal hardware, or severe spine arthritis so they scan the best bone sites.
\end{itemize}

\section{Preparing for the Scan}
\begin{itemize}
\item \textbf{Skip Calcium Supplements for 24 Hours:} Do not take calcium pills or multivitamins containing calcium for at least 24 hours before your scan. Undissolved calcium tablets in your stomach or intestines can show up over your spine and falsely raise your score.
\item \textbf{Eat Normally and Wear Metal-Free Clothes:} You can eat normally and take your other regular medicines as usual. Wear loose, comfortable clothes with no metal zippers, buttons, snaps, underwire bras, or belt buckles over your hips and lower back (or change into a gown) and empty your pockets.
\item \textbf{Wait After Recent Contrast or Nuclear Scans:} Call the clinic to reschedule (usually waiting 7--14 days) if you recently had a barium exam, contrast CT scan, or nuclear medicine test, since leftover dye or tracers skew DXA readings.
\item \textbf{Share Key Details and Past Results:} Tell the technologist if you might be pregnant or have spine or hip metal hardware. Bring past DXA reports if you have them, and try to return to the same imaging center and machine for repeat scans.
\end{itemize}

\section{During Scan}
\begin{itemize}
\item \textbf{Comfortable, Open-Air Table:} You lie flat on your back on a wide, open padded table---there is no enclosed tunnel. For the lower-back scan, a soft foam box supports your lower legs to flatten your spine; for the hip scan, your foot rests gently in a triangular holder that turns your leg slightly inward.
\item \textbf{Relaxed Breathing and Staying Still:} Breathe normally and quietly, avoiding talking or shifting your hips or legs while the overhead arm glides above you; the technologist may occasionally ask you to hold your breath for a few seconds to prevent blurring.
\item \textbf{Completely Painless:} You will not feel the X-rays at all. The overhead arm moves slowly over your body without touching you and makes only a quiet, gentle hum.
\item \textbf{Tell the Staff If Your Back or Hip Aches:} Let the technologist know if lying flat feels uncomfortable so they can adjust pillows or supports safely.
\item \textbf{Fast Visit:} Expect only 10--20 minutes total on the table, with each spine, hip, or forearm scan taking just 1--3 minutes.
\end{itemize}

\section{Normal Appearance}
\begin{itemize}
\item \textbf{Healthy Bone-Density Scores:} For postmenopausal women and men aged 50 or older, a T-score of -1.0 or higher at the lower spine, hip neck, and total hip is considered normal. For younger adults and children, a Z-score above -2.0 is within the expected range for age.
\item \textbf{Straight, Even Lower Spine:} The four lower back bones (L1--L4) line up straight and centered without metal objects over them, and bone density normally rises gently from the top bone (L1) down to L4.
\item \textbf{Proper Hip Alignment:} The thigh bone lies straight and slightly turned inward so the hip neck and total hip zones can be measured accurately.
\item \textbf{Healthy Spine Shapes on Side Views:} If a lateral spine fracture check (VFA) is added, healthy back bones look like neat, even rectangles with no wedging or flattening.
\end{itemize}

\section{Observing for Abnormal Patterns}
\begin{itemize}
\item \textbf{Low Bone Mass (Osteopenia) and Osteoporosis:} For postmenopausal women and men 50 and older, a lowest T-score between -1.0 and -2.5 shows osteopenia, and -2.5 or lower confirms osteoporosis. In younger adults and children, a Z-score of -2.0 or lower signals bone density below the expected range for age.
\item \textbf{One Spine Bone Looking Much Denser Than Its Neighbors:} If one lower-back vertebra has a T-score more than 1.0 point higher than the bone next to it, that usually means arthritis spurs or a collapsed fracture is skewing that level. Per guidelines, doctors exclude that skewed bone and average the remaining healthy ones (at least two).
\item \textbf{Silent Spine Compression Fractures:} On a side-view VFA scan, a vertebra that has lost 20\% or more of its height---looking wedged, pinched in the middle, or crushed---confirms a fragility fracture that signals high risk for future breaks even if standard bone density is only in the osteopenia range.
\item \textbf{Extra Calcium and Arthritis Overlying the Spine:} Heavy arthritis bone spurs, calcium in the abdominal aorta, or rare bone conditions can show up as bright dense spots that falsely inflate spine density and hide true osteoporosis.
\item \textbf{Positioning or Foreign-Object Quirks:} If the hip is not turned inward enough, hip density can look falsely higher. Leftover barium, contrast dye, metal clothing fasteners, surgical hardware, or undissolved calcium tablets can also create bright artifacts that throw off regional scores.
\end{itemize}

\section{Dealing with Uncertainties}
\begin{itemize}
\item \textbf{Arthritis, Calcium, and Supplements Skewing Readings:} Heavy spine arthritis, old fractures, calcified blood vessels, metal hardware, or leftover contrast/calcium pills can make spine bone density look artificially strong even when bones are actually fragile. Similarly, taking strontium supplements causes DXA to overestimate bone density because strontium absorbs X-rays more strongly than calcium.
\item \textbf{Positioning and Different-Machine Differences:} Slight shifts in hip angle, curved spines (scoliosis), or comparing scans taken on two different scanner brands can look like fake bone loss or fake improvement.
\item \textbf{Using Backup Bone Sites:} When individual spine bones are skewed by arthritis, doctors exclude them from the average (keeping at least two valid levels). If fewer than two spine levels or neither hip can be measured reliably, the non-dominant forearm (33\% radius) provides a trustworthy alternative.
\item \textbf{Fractures Despite ``Borderline'' Scores:} Because standard DXA does not show fine inner bone architecture, some people with osteopenia still break bones easily. Doctors resolve this by adding a FRAX risk score, Trabecular Bone Score (TBS), or lateral spine fracture check (VFA).
\item \textbf{Repeating or Switching Tests When Needed:} If contrast dye or positioning caused an error, the scan can be repeated once cleared (after 7--14 days). In severe spinal arthritis, specialized quantitative CT (QCT) or a visit with a bone-health specialist can give clearer answers.
\end{itemize}

\section{Augmenting Methods}
\begin{itemize}
\item \textbf{Trabecular Bone Score (TBS) and Spine Fracture Check (VFA):} TBS software analyzes the fine bone texture on your existing spine DXA image, and VFA checks for silent spine fractures---both sharpening fracture-risk estimates without extra fuss.
\item \textbf{FRAX Fracture Risk Calculator:} Combines your hip bone density (and optional TBS) with personal risk factors like age, past breaks, smoking, or steroid use to calculate your 10-year odds of a major fracture.
\item \textbf{Bone-Health Blood Tests:} Simple labs---calcium, vitamin D, parathyroid hormone, kidney function, thyroid hormone, blood count, and optional bone-turnover markers---spot underlying causes of bone thinning and track early response to treatment.
\item \textbf{Quantitative CT (QCT) and Spine X-Rays or MRI:} QCT measures true 3D inner bone density unaffected by spine arthritis or artery calcium, while standard spine X-rays or MRI confirm painful new fractures and rule out other bone problems.
\item \textbf{AI Opportunistic Screening on Routine CT Scans:} Newer AI tools can automatically check spine bone strength and spot hidden fractures on routine chest or belly CT scans done for other reasons, flagging people who should get a formal DXA scan.
\end{itemize}

\section{Comparison with Alternative Modalities}
\begin{itemize}
\item \textbf{Versus Quantitative CT (QCT):} QCT measures true 3D inner bone density without interference from arthritis or calcified blood vessels. However, DXA is the global diagnostic gold standard with dramatically lower radiation, lower cost, and better repeatability for tracking treatment over time.
\item \textbf{Versus Heel Ultrasound (QUS):} A portable heel ultrasound uses zero radiation and can offer a quick preliminary risk check, but it cannot officially diagnose osteoporosis using WHO T-scores or track medication response like central DXA.
\item \textbf{Versus Peripheral DEXA (p-DEXA):} Compact, box-like p-DEXA scanners (sometimes available in clinics, pharmacies, or mobile health vans) measure the wrist, finger, forearm, or heel in a few minutes to flag possible bone loss, but they provide less detail than a central DXA of the hip and lower spine and usually require central DXA follow-up if risk appears elevated.
\item \textbf{Versus Standard Spine X-Rays and MRI:} Standard X-rays and MRI are best for diagnosing a sudden, painful broken vertebra or checking bone marrow swelling, whereas DXA measures overall bone strength and screens for quiet, painless fractures via VFA.
\item \textbf{Versus Nuclear Medicine Bone Scan:} Despite their similar names, a DXA ``bone density scan'' measures mineral strength to check for osteopenia and osteoporosis, whereas a nuclear ``bone scan'' uses an injected radioactive tracer to highlight areas of active bone turnover or healing from infections, fractures, or tumors.
\end{itemize}

\section{Follow-up Steps}
\begin{itemize}
\item \textbf{Right After the Exam:} DXA has zero recovery time---you can get dressed, resume your normal day immediately, and restart your calcium supplements right away.
\item \textbf{Clear Report and Fracture Risk Estimate:} Your doctor receives your bone density numbers, T-scores or Z-scores, any notes on excluded arthritis-affected spine levels, and your 10-year FRAX fracture risk to guide whether lifestyle changes, vitamin D/calcium, or osteoporosis medications make sense.
\item \textbf{Checking for Underlying Causes:} If osteoporosis or fragility fractures are newly found, your doctor will typically order basic metabolic blood tests (calcium, vitamin D, kidney, blood count, thyroid), check spine X-rays if you have lost height, and consider a bone-health or endocrinology specialist for severe cases.
\item \textbf{When to Repeat Your Scan:} Ideally returning to the same scanner for accurate comparison, DXA is usually repeated every 1--2 years after starting or changing bone-strengthening medication, and less often (every 3--5 years or longer) once bone density is stable or in low-risk osteopenia.
\end{itemize}

\section{Key References}
\begin{itemize}
\item International Society for Clinical Densitometry (ISCD). \textit{2023 ISCD Official Positions --- Adult and Pediatric}. Middletown, CT: ISCD; 2023.
\item American College of Radiology (ACR). \textit{ACR-SPR-SSR Practice Parameter for the Performance of Dual-Energy X-Ray Absorptiometry (DXA)} and \textit{ACR Appropriateness Criteria: Osteoporosis and Bone Mineral Density}. Reston, VA: ACR; updated 2022.
\item LeBoff MS, Greenspan SL, Insogna KL, et al. The clinician's guide to prevention and treatment of osteoporosis (Bone Health and Osteoporosis Foundation). \textit{Osteoporosis International}. 2022;33(10):2049--2102.
\item World Health Organization (WHO). \textit{Assessment of Fracture Risk and Its Application to Screening for Postmenopausal Osteoporosis}. WHO Technical Report Series 843. Geneva: WHO; 1994.
\item U.S. National Library of Medicine. \textit{Bone Density Scan}. MedlinePlus, National Institutes of Health (NIH). Available at: \url{https://medlineplus.gov/lab-tests/bone-density-scan/}.
\end{itemize}

\section{Further Reading}
\begin{itemize}
  \item Wikipedia contributors. \textit{Dual-energy X-ray absorptiometry}. Wikipedia, The Free Encyclopedia. Available at: \url{https://en.wikipedia.org/wiki/Dual-energy_X-ray_absorptiometry}.
\end{itemize}
